
\documentclass{article}
\usepackage{colortbl}
\usepackage{makecell}
\usepackage{multirow}
\usepackage{supertabular}

\begin{document}

\newcounter{utterance}

\centering \large Interaction Transcript for game `wordle', experiment `medium\_frequency\_words\_no\_clue\_no\_critic', episode 9 with mistral{-}large{-}3{-}azure.
\vspace{24pt}

{ \footnotesize  \setcounter{utterance}{1}
\setlength{\tabcolsep}{0pt}
\begin{supertabular}{c@{$\;$}|p{.15\linewidth}@{}p{.15\linewidth}p{.15\linewidth}p{.15\linewidth}p{.15\linewidth}p{.15\linewidth}}
    \# & \multicolumn{2}{c}{Player} && \multicolumn{2}{c}{Game Master} \\
    \hline

    \theutterance \stepcounter{utterance}  
    & & & \multicolumn{4}{p{0.6\linewidth}}{
        \cellcolor[rgb]{0.9,0.9,0.9}{
            \makecell[{{p{\linewidth}}}]{
                \texttt{\tiny{[P1$\langle$GM]}}
                \texttt{Вы — языковой волшебник, который любит угадывать слова, используя заданные правила.} \\
\\ 
\texttt{Добро пожаловать в Wordle! У вас есть шесть попыток, чтобы угадать целевое слово — допустимое слово из пяти букв. Пожалуйста, используйте теги "объяснение:" и "догадка:", чтобы дать краткое объяснение для каждой попытки.} \\
\\ 
\texttt{Например, если ваша догадка — "слово", ваш ответ должен быть} \\
\texttt{объяснение: это распространённое слово из пяти букв, и я начинаю угадывать с этого слова.} \\
\texttt{догадка: слово} \\
\\ 
\texttt{После каждой попытки ваш ответ будет проверен, и вы получите обратную связь, показывающую, какие буквы верны (зелёный), какие буквы верны, но стоят не на своём месте (жёлтый), и какие буквы неверны (красный). Эта обратная связь может быть полезна для определения того, какие буквы включить или исключить в следующей попытке.} \\
\\ 
\texttt{Например, обратная связь для "слово" может быть:} \\
\texttt{оценка\_догадки: с<жёлтый> л<жёлтый> о<зелёный> в<жёлтый> о<красный>} \\
\\ 
\texttt{В explanation должны быть детали о том, как оценка\_догадки используется, чтобы прийти к новой догадке.} \\
\\ 
\texttt{Не генерируйте никакого другого текста.} \\
\\ 
\texttt{Начнём с вашей первой догадки.} \\
            }
        }
    }
     \\ \\

    \theutterance \stepcounter{utterance}  
    & \multicolumn{4}{p{0.6\linewidth}}{
        \cellcolor[rgb]{0.561,0.737,0.561}{
            \makecell[{{p{\linewidth}}}]{
                \texttt{\tiny{[P1$\rangle$GM]}}
                \texttt{объяснение: Начну с распространённого слова из пяти букв, содержащего часто встречающиеся гласные (О, А, Е) и согласные (Т, Р, Н, С, К). Это поможет быстро проверить популярные буквы.} \\
\texttt{догадка: ротан} \\
            }
        }
    }
    & & \\ \\

    \theutterance \stepcounter{utterance}  
    & & & \multicolumn{2}{p{0.3\linewidth}}{
        \cellcolor[rgb]{0.9,0.9,0.9}{
            \makecell[{{p{\linewidth}}}]{
                \texttt{\tiny{[GM$|$GM]}}
                \texttt{Error: The guessed word is not a valid word for this game.} \\
            }
        }
    }
    & & \\ \\

    \theutterance \stepcounter{utterance}  
    & & & \multicolumn{4}{p{0.6\linewidth}}{
        \cellcolor[rgb]{0.9,0.9,0.9}{
            \makecell[{{p{\linewidth}}}]{
                \texttt{\tiny{[P1$\langle$GM]}}
                \texttt{Ваша догадка не является допустимым словом для этой игры. Пожалуйста, попробуйте ещё раз.} \\
\\ 
\texttt{Дайте ответ только в этом формате.} \\
\texttt{объяснение: подробности} \\
\texttt{догадка: слово} \\
            }
        }
    }
     \\ \\

    \theutterance \stepcounter{utterance}  
    & \multicolumn{4}{p{0.6\linewidth}}{
        \cellcolor[rgb]{0.561,0.737,0.561}{
            \makecell[{{p{\linewidth}}}]{
                \texttt{\tiny{[P1$\rangle$GM]}}
                \texttt{объяснение: Выберу распространённое слово из пяти букв с часто встречающимися буквами (А, Р, О, С, Т), чтобы проверить их наличие и позиции.} \\
\texttt{догадка: арост} \\
            }
        }
    }
    & & \\ \\

    \theutterance \stepcounter{utterance}  
    & & & \multicolumn{2}{p{0.3\linewidth}}{
        \cellcolor[rgb]{0.9,0.9,0.9}{
            \makecell[{{p{\linewidth}}}]{
                \texttt{\tiny{[GM$|$GM]}}
                \texttt{Error: The guessed word is not a valid word for this game.} \\
            }
        }
    }
    & & \\ \\

    \theutterance \stepcounter{utterance}  
    & & & \multicolumn{4}{p{0.6\linewidth}}{
        \cellcolor[rgb]{0.9,0.9,0.9}{
            \makecell[{{p{\linewidth}}}]{
                \texttt{\tiny{[P1$\langle$GM]}}
                \texttt{Ваша догадка не является допустимым словом для этой игры. Пожалуйста, попробуйте ещё раз.} \\
\\ 
\texttt{Дайте ответ только в этом формате.} \\
\texttt{объяснение: подробности} \\
\texttt{догадка: слово} \\
            }
        }
    }
     \\ \\

    \theutterance \stepcounter{utterance}  
    & \multicolumn{4}{p{0.6\linewidth}}{
        \cellcolor[rgb]{0.561,0.737,0.561}{
            \makecell[{{p{\linewidth}}}]{
                \texttt{\tiny{[P1$\rangle$GM]}}
                \texttt{объяснение: Начну с классического стартового слова в Wordle, которое содержит распространённые гласные (А, Е) и согласные (К, Н, Т). Это поможет быстро сузить круг возможных букв.} \\
\texttt{догадка: танец} \\
            }
        }
    }
    & & \\ \\

    \theutterance \stepcounter{utterance}  
    & & & \multicolumn{2}{p{0.3\linewidth}}{
        \cellcolor[rgb]{0.9,0.9,0.9}{
            \makecell[{{p{\linewidth}}}]{
                \texttt{\tiny{[GM$|$GM]}}
                \texttt{attempts = 1} \\
\texttt{target\_word = отдых} \\
\texttt{guess = танец} \\
\texttt{guess\_feedback = т<yellow> а<red> н<red> е<red> ц<red>} \\
            }
        }
    }
    & & \\ \\

    \theutterance \stepcounter{utterance}  
    & & & \multicolumn{4}{p{0.6\linewidth}}{
        \cellcolor[rgb]{0.9,0.9,0.9}{
            \makecell[{{p{\linewidth}}}]{
                \texttt{\tiny{[P1$\langle$GM]}}
                \texttt{оценка\_догадки: т<жёлтый> а<красный> н<красный> е<красный> ц<красный>} \\
\\ 
\texttt{Дайте ответ только в этом формате.} \\
\texttt{объяснение: подробности} \\
\texttt{догадка: слово} \\
            }
        }
    }
     \\ \\

    \theutterance \stepcounter{utterance}  
    & \multicolumn{4}{p{0.6\linewidth}}{
        \cellcolor[rgb]{0.561,0.737,0.561}{
            \makecell[{{p{\linewidth}}}]{
                \texttt{\tiny{[P1$\rangle$GM]}}
                \texttt{объяснение: Буква **Т** присутствует в слове, но не на первом месте (жёлтый). Все остальные буквы (**А, Н, Е, Ц**) отсутствуют (красный). Исключаю их и выбираю новое слово с **Т** в другой позиции, а также добавляю новые распространённые буквы (**Р, О, С, Л**).} \\
\texttt{догадка: торос} \\
            }
        }
    }
    & & \\ \\

    \theutterance \stepcounter{utterance}  
    & & & \multicolumn{2}{p{0.3\linewidth}}{
        \cellcolor[rgb]{0.9,0.9,0.9}{
            \makecell[{{p{\linewidth}}}]{
                \texttt{\tiny{[GM$|$GM]}}
                \texttt{attempts = 2} \\
\texttt{target\_word = отдых} \\
\texttt{guess = торос} \\
\texttt{guess\_feedback = т<yellow> о<yellow> р<red> о<red> с<red>} \\
            }
        }
    }
    & & \\ \\

    \theutterance \stepcounter{utterance}  
    & & & \multicolumn{4}{p{0.6\linewidth}}{
        \cellcolor[rgb]{0.9,0.9,0.9}{
            \makecell[{{p{\linewidth}}}]{
                \texttt{\tiny{[P1$\langle$GM]}}
                \texttt{оценка\_догадки: т<жёлтый> о<жёлтый> р<красный> о<красный> с<красный>} \\
\\ 
\texttt{Дайте ответ только в этом формате.} \\
\texttt{объяснение: подробности} \\
\texttt{догадка: слово} \\
            }
        }
    }
     \\ \\

    \theutterance \stepcounter{utterance}  
    & \multicolumn{4}{p{0.6\linewidth}}{
        \cellcolor[rgb]{0.561,0.737,0.561}{
            \makecell[{{p{\linewidth}}}]{
                \texttt{\tiny{[P1$\rangle$GM]}}
                \texttt{объяснение:} \\
\texttt{{-} Буква **Т** снова жёлтая → не на первом месте, но есть в слове.} \\
\texttt{{-} **О** тоже жёлтая → присутствует, но не на второй позиции (вторая **О** красная → значит, в слове только одна **О**).} \\
\texttt{{-} Буквы **Р, С** отсутствуют (красный).} \\
\texttt{{-} Исключаю **Р, С** и ищу слово с **Т** и одной **О**, добавляя новые буквы (**К, Л, Д, П**).} \\
\texttt{догадка: плоть} \\
            }
        }
    }
    & & \\ \\

    \theutterance \stepcounter{utterance}  
    & & & \multicolumn{2}{p{0.3\linewidth}}{
        \cellcolor[rgb]{0.9,0.9,0.9}{
            \makecell[{{p{\linewidth}}}]{
                \texttt{\tiny{[GM$|$GM]}}
                \texttt{attempts = 3} \\
\texttt{target\_word = отдых} \\
\texttt{guess = плоть} \\
\texttt{guess\_feedback = п<red> л<red> о<yellow> т<yellow> ь<red>} \\
            }
        }
    }
    & & \\ \\

    \theutterance \stepcounter{utterance}  
    & & & \multicolumn{4}{p{0.6\linewidth}}{
        \cellcolor[rgb]{0.9,0.9,0.9}{
            \makecell[{{p{\linewidth}}}]{
                \texttt{\tiny{[P1$\langle$GM]}}
                \texttt{оценка\_догадки: п<красный> л<красный> о<жёлтый> т<жёлтый> ь<красный>} \\
\\ 
\texttt{Дайте ответ только в этом формате.} \\
\texttt{объяснение: подробности} \\
\texttt{догадка: слово} \\
            }
        }
    }
     \\ \\

    \theutterance \stepcounter{utterance}  
    & \multicolumn{4}{p{0.6\linewidth}}{
        \cellcolor[rgb]{0.561,0.737,0.561}{
            \makecell[{{p{\linewidth}}}]{
                \texttt{\tiny{[P1$\rangle$GM]}}
                \texttt{объяснение:} \\
\texttt{{-} **О** и **Т** снова жёлтые → они есть в слове, но не на 3{-}й и 4{-}й позициях соответственно.} \\
\texttt{{-} **П, Л, Ь** отсутствуют (красный) → исключаю их.} \\
\texttt{{-} Нужно слово с **О** и **Т**, но без **П, Л, Р, С, А, Н, Е, Ц** (из предыдущих попыток).} \\
\texttt{{-} Проверяю **Д, К, М, И** как новые буквы, сохраняя **Т** и **О** в других позициях.} \\
\texttt{догадка: мотик} \\
            }
        }
    }
    & & \\ \\

    \theutterance \stepcounter{utterance}  
    & & & \multicolumn{2}{p{0.3\linewidth}}{
        \cellcolor[rgb]{0.9,0.9,0.9}{
            \makecell[{{p{\linewidth}}}]{
                \texttt{\tiny{[GM$|$GM]}}
                \texttt{attempts = 4} \\
\texttt{target\_word = отдых} \\
\texttt{guess = мотик} \\
\texttt{guess\_feedback = м<red> о<yellow> т<yellow> и<red> к<red>} \\
            }
        }
    }
    & & \\ \\

    \theutterance \stepcounter{utterance}  
    & & & \multicolumn{4}{p{0.6\linewidth}}{
        \cellcolor[rgb]{0.9,0.9,0.9}{
            \makecell[{{p{\linewidth}}}]{
                \texttt{\tiny{[P1$\langle$GM]}}
                \texttt{оценка\_догадки: м<красный> о<жёлтый> т<жёлтый> и<красный> к<красный>} \\
\\ 
\texttt{Дайте ответ только в этом формате.} \\
\texttt{объяснение: подробности} \\
\texttt{догадка: слово} \\
            }
        }
    }
     \\ \\

    \theutterance \stepcounter{utterance}  
    & \multicolumn{4}{p{0.6\linewidth}}{
        \cellcolor[rgb]{0.561,0.737,0.561}{
            \makecell[{{p{\linewidth}}}]{
                \texttt{\tiny{[P1$\rangle$GM]}}
                \texttt{объяснение:} \\
\texttt{{-} **О** и **Т** снова жёлтые → они есть, но не на 2{-}й и 3{-}й позициях.} \\
\texttt{{-} Исключаю **М, И, К** (красный) и все ранее проверенные буквы (**П, Л, Р, С, А, Н, Е, Ц, Ь**).} \\
\texttt{{-} Остаются распространённые буквы: **Д, В, Г, Б, У, Ы**.} \\
\texttt{{-} Проверяю слово с **Т** в конце (возможная позиция) и **О** в начале или середине, добавляя новую букву **Д**.} \\
\texttt{догадка: доток} \\
            }
        }
    }
    & & \\ \\

    \theutterance \stepcounter{utterance}  
    & & & \multicolumn{2}{p{0.3\linewidth}}{
        \cellcolor[rgb]{0.9,0.9,0.9}{
            \makecell[{{p{\linewidth}}}]{
                \texttt{\tiny{[GM$|$GM]}}
                \texttt{Error: The guessed word is not a valid word for this game.} \\
            }
        }
    }
    & & \\ \\

    \theutterance \stepcounter{utterance}  
    & & & \multicolumn{4}{p{0.6\linewidth}}{
        \cellcolor[rgb]{0.9,0.9,0.9}{
            \makecell[{{p{\linewidth}}}]{
                \texttt{\tiny{[P1$\langle$GM]}}
                \texttt{Ваша догадка не является допустимым словом для этой игры. Пожалуйста, попробуйте ещё раз.} \\
\\ 
\texttt{Дайте ответ только в этом формате.} \\
\texttt{объяснение: подробности} \\
\texttt{догадка: слово} \\
            }
        }
    }
     \\ \\

    \theutterance \stepcounter{utterance}  
    & \multicolumn{4}{p{0.6\linewidth}}{
        \cellcolor[rgb]{0.561,0.737,0.561}{
            \makecell[{{p{\linewidth}}}]{
                \texttt{\tiny{[P1$\rangle$GM]}}
                \texttt{объяснение:} \\
\texttt{{-} **О** и **Т** должны быть в слове, но не на 2{-}й и 3{-}й позициях (из предыдущих попыток).} \\
\texttt{{-} Исключены буквы: **А, Н, Е, Ц, Р, С, П, Л, Ь, М, И, К**.} \\
\texttt{{-} Остаются варианты с **Д, В, Г, Б, У, Ы** + возможные **О** и **Т** в других местах.} \\
\texttt{{-} Проверяю слово с **Т** на 4{-}й позиции и **О** на 1{-}й, добавляя **Д** и новую букву **В**.} \\
\texttt{догадка: довод} \\
            }
        }
    }
    & & \\ \\

    \theutterance \stepcounter{utterance}  
    & & & \multicolumn{2}{p{0.3\linewidth}}{
        \cellcolor[rgb]{0.9,0.9,0.9}{
            \makecell[{{p{\linewidth}}}]{
                \texttt{\tiny{[GM$|$GM]}}
                \texttt{attempts = 5} \\
\texttt{target\_word = отдых} \\
\texttt{guess = довод} \\
\texttt{guess\_feedback = д<yellow> о<yellow> в<red> о<red> д<red>} \\
            }
        }
    }
    & & \\ \\

    \theutterance \stepcounter{utterance}  
    & & & \multicolumn{4}{p{0.6\linewidth}}{
        \cellcolor[rgb]{0.9,0.9,0.9}{
            \makecell[{{p{\linewidth}}}]{
                \texttt{\tiny{[P1$\langle$GM]}}
                \texttt{оценка\_догадки: д<жёлтый> о<жёлтый> в<красный> о<красный> д<красный>} \\
\\ 
\texttt{Дайте ответ только в этом формате.} \\
\texttt{объяснение: подробности} \\
\texttt{догадка: слово} \\
            }
        }
    }
     \\ \\

    \theutterance \stepcounter{utterance}  
    & \multicolumn{4}{p{0.6\linewidth}}{
        \cellcolor[rgb]{0.561,0.737,0.561}{
            \makecell[{{p{\linewidth}}}]{
                \texttt{\tiny{[P1$\rangle$GM]}}
                \texttt{объяснение:} \\
\texttt{{-} **Д** и **О** жёлтые → есть в слове, но не на 1{-}й и 2{-}й позициях (вторая **О** и второй **Д** красные → только одна **О** и одна **Д**).} \\
\texttt{{-} **В** отсутствует (красный) → исключаю.} \\
\texttt{{-} Остаются буквы: **Г, Б, У, Ы, Т** (из предыдущих попыток **Т** должна быть в слове).} \\
\texttt{{-} Слово должно содержать **Т**, **Д**, **О**, но не на проверенных позициях.} \\
\texttt{{-} Проверяю **Т** на 5{-}й позиции, **О** на 3{-}й, **Д** на 4{-}й, добавляя новую букву **Г**.} \\
\texttt{догадка: гогот} \\
            }
        }
    }
    & & \\ \\

    \theutterance \stepcounter{utterance}  
    & & & \multicolumn{2}{p{0.3\linewidth}}{
        \cellcolor[rgb]{0.9,0.9,0.9}{
            \makecell[{{p{\linewidth}}}]{
                \texttt{\tiny{[GM$|$GM]}}
                \texttt{attempts = 6} \\
\texttt{target\_word = отдых} \\
\texttt{guess = гогот} \\
\texttt{guess\_feedback = г<red> о<yellow> г<red> о<red> т<yellow>} \\
            }
        }
    }
    & & \\ \\

    \theutterance \stepcounter{utterance}  
    & & & \multicolumn{2}{p{0.3\linewidth}}{
        \cellcolor[rgb]{0.9,0.9,0.9}{
            \makecell[{{p{\linewidth}}}]{
                \texttt{\tiny{[GM$|$GM]}}
                \texttt{game\_result = LOSS} \\
            }
        }
    }
    & & \\ \\

\end{supertabular}
}

\end{document}
